\documentclass[11pt]{article}


\usepackage{fullpage}
\usepackage{graphicx}
\usepackage{amsmath}
\usepackage{amssymb}
\usepackage{amsthm}
\usepackage{fancyvrb}

\newcommand{\myname}{Mehshan Mustafa}

\newenvironment{theorem}[2][Theorem]{\begin{trivlist}
\item[\hskip \labelsep {\bfseries #1}\hskip \labelsep {\bfseries #2.}]}{\end{trivlist}}
\newenvironment{lemma}[2][Lemma]{\begin{trivlist}
\item[\hskip \labelsep {\bfseries #1}\hskip \labelsep {\bfseries #2.}]}{\end{trivlist}}
\newenvironment{exercise}[2][Exercise]{\begin{trivlist}
\item[\hskip \labelsep {\bfseries #1}\hskip \labelsep {\bfseries #2.}]}{\end{trivlist}}
\newenvironment{problem}[2][Problem]{\begin{trivlist}
\item[\hskip \labelsep {\bfseries #1}\hskip \labelsep {\bfseries #2.}]}{\end{trivlist}}
\newenvironment{question}[2][Question]{\begin{trivlist}
\item[\hskip \labelsep {\bfseries #1}\hskip \labelsep {\bfseries #2.}]}{\end{trivlist}}
\newenvironment{corollary}[2][Corollary]{\begin{trivlist}
\item[\hskip \labelsep {\bfseries #1}\hskip \labelsep {\bfseries #2.}]}{\end{trivlist}}
\newenvironment{solution}{\begin{proof}[Solution]}{\end{proof}}
\newenvironment{idea}[2][Proof Idea.]{\textit{#1} #2}



\parindent0in
\pagestyle{plain}
\thispagestyle{plain}

\usepackage{csquotes}
\usepackage[shortlabels]{enumitem}

\newcommand{\dated}{\today}
\newcommand{\token}[1]{\langle \text{#1} \rangle}

\begin{document}

\textbf{Introduction to the Theory of
Computation}\hfill\textbf{\myname}\\[0.01in]
\textbf{Chapter 10: Advanced Topics in Complexity Theory}\hfill\textbf{\dated}\\
\smallskip\hrule\bigskip

\begin{problem}{10.18}
Prove that if $A \subseteq \{0, 1\}^*$, a family of branching programs $(B_1,B_2, . . .)$ exists wherein each $B_n$ accepts exactly the strings in $A$ of length $n$ and is bounded in size by a constant times $n$.
\end{problem}

\begin{proof}
$A$ is a regular language. Therefore, a DFA $D$ exists that accepts every string in $A$.    Construct each $B_n$ as follows.
\begin{enumerate}
\item States of $D$ are query nodes in $B_n$, and transitions of $D$ are outgoing edges in $B_n$.
\item Each input symbol $w_i$ in string $w_1w_2 \cdots w_n$ corresponds to variable $x_i$ in $B_n$.
\item $B_n$ includes all nodes and transitions of $D$ required for accepting all strings of length $n$ and loops are unrolled as many times as needed.
\item Start node in $D$ is also the start node in $B_n$.
\item Each query node $q$ in $B_n$ is labeled $x_{p + 1}$, where $p$ is the length of path from start node to $q$.
\item Finally, accepts node of $D$ are replaced by single output node labeled 1, and all reject paths end at output node labeled 0.     
\end{enumerate}

\end{proof}

\end{document}